% Source: Gerzensee "Real Analysis" slides 26-48 (open balls, bounded sets,
% interior/closure, open and closed sets, topology axioms, the sigma-algebra
% remark, sequences, subsequences, limits, monotone subsequences,
% Bolzano-Weierstrass, coordinatewise convergence).
% Supplementary: Fukuda (2026), Sec. 7.2-7.3 (pp. 214-238) and Sec. 8.1
% (Compact Sets, pp. 261-270).
% External references actually cited in this chapter: Aliprantis-Border
% (2006), Rudin (1976).

\chapter[Topology of Metric Spaces]{Topology of Metric Spaces: Open Sets, Sequences, and Compactness}
\label{ch:topology}

\section{Motivation: What ``Nearby'' Has to Mean}\label{sec:top-motivation}

Three of the most-used theorems in economic theory --- that a continuous function
on a compact set attains its maximum, that a bounded sequence has a convergent
subsequence, and that a contraction on a complete space has a unique fixed point
--- are statements about which sets contain their limit points and which
sequences converge. This chapter builds that vocabulary.

The order of business follows the slides
\autocite[slides 26--48]{fukuda2026analysis}: open balls, then interior and
closure, then the three axioms that isolate what a topology is, then sequences
and their limits, then the two results that make compactness usable ---
Bolzano--Weierstrass and Heine--Borel. Throughout, $(X,d)$ is a metric space; the
specialisation to $\R^n$ is stated whenever it is genuinely used, because several
of the results below are false in general metric spaces and knowing which is the
point of stating them abstractly.

\section{Open and Closed Sets}\label{sec:top-openclosed}

\begin{definition}[Interior, closure, boundary]\label{def:interior-closure}
	Let $(X,d)$ be a metric space and $A\subseteq X$.
	\begin{enumerate}[(i)]
		\item A point $a\in X$ is \dfnas{interior point}{interior to} $A$ if
		      $\ball{\varepsilon}{a}\subseteq A$ for some $\varepsilon>0$. The
		      \dfn{interior} $\Int A$ is the set of interior points of $A$; by
		      definition $\Int A\subseteq A$. The set $A$ is \dfnas{open set}{open} if
		      $A=\Int A$.
		\item A point $a\in X$ is a \dfnas{point of closure}{point of closure} of $A$
		      if $A\cap\ball{\varepsilon}{a}\neq\emptyset$ for every $\varepsilon>0$.
		      The \dfn{closure} $\closure{A}$ is the set of points of closure; by
		      definition $A\subseteq\closure{A}$. The set $A$ is \dfnas{closed
			      set}{closed} if $A=\closure{A}$.
		\item The \dfn{boundary} is $\partial A\eqdef\closure{A}\setminus\Int A$.
	\end{enumerate}
	\nidx{$\Int A$}{$\Int A$, interior}\nidx{$\overline{A}$}{$\closure{A}$, closure}
\end{definition}

\begin{eg}[Closure can add points that are ``reached but not attained'']
	\label{eg:closure-1overn}
	In $\R$ with $d(x,y)=\abs{x-y}$ let $A=\{1/n : n\in\N\}$. Then $0\in\closure A$:
	for any $\varepsilon>0$, \autoref{prop:archimedean} supplies $n$ with
	$1/n<\varepsilon$, so $1/n\in A\cap\ball{\varepsilon}{0}$. Since $0\notin A$,
	$A$ is not closed, and $\closure A=A\cup\{0\}$. Also $\Int A=\emptyset$, since
	every ball around $1/n$ contains irrational numbers. This is Example~7.1 of
	\textcite{fukuda2026notes}, illustrated on
	\autocite[slide 33]{fukuda2026analysis}.
\end{eg}

\begin{proposition}[Duality of interior and closure]\label{prop:interior-closure-duality}
	For every $A\subseteq X$,
	\[
		\comp{\bigl(\closure{A}\bigr)}=\Int\bigl(\comp{A}\bigr),
		\qquad\text{equivalently}\qquad
		\comp{\bigl(\Int A\bigr)}=\closure{\bigl(\comp{A}\bigr)} .
	\]
	Consequently $A$ is open if and only if $\comp{A}$ is closed, and $A$ is closed
	if and only if $\comp{A}$ is open.
\end{proposition}

\begin{proof}
	A point $a$ fails to be a point of closure of $A$ if and only if some
	$\varepsilon>0$ has $A\cap\ball{\varepsilon}{a}=\emptyset$, that is
	$\ball{\varepsilon}{a}\subseteq\comp{A}$, that is $a\in\Int(\comp{A})$. This is
	the first identity; the second follows by applying the first to $\comp{A}$ and
	taking complements.

	For the consequence: $A$ open means $A=\Int A$, hence
	$\comp{A}=\comp{(\Int A)}=\closure{(\comp{A})}$, which is $\comp A$ closed; and
	conversely.
\end{proof}

This is Proposition~7.5 and Corollary~7.1 of \textcite{fukuda2026notes}
\autocite[slide 35]{fukuda2026analysis}. Almost every proof about closed sets
below is obtained by complementing a proof about open sets.

\begin{proposition}[Open balls are open]\label{prop:ball-open}
	In any metric space, $\ball{\lambda}{a}$ is open.
\end{proposition}

\begin{proof}
	Let $x\in\ball{\lambda}{a}$ and put $\rho\eqdef d(a,x)<\lambda$, so
	$\lambda-\rho>0$. For $y\in\ball{\lambda-\rho}{x}$ the triangle inequality gives
	\[
		d(a,y)\leq d(a,x)+d(x,y)<\rho+(\lambda-\rho)=\lambda,
	\]
	so $y\in\ball{\lambda}{a}$. Hence $\ball{\lambda-\rho}{x}\subseteq
		\ball{\lambda}{a}$ and $x\in\Int\ball{\lambda}{a}$. As $x$ was arbitrary,
	$\ball{\lambda}{a}=\Int\ball{\lambda}{a}$.
\end{proof}

\begin{proposition}[$\Int A$ is open and $\closure{A}$ is closed]\label{prop:int-open}
	For every $A\subseteq X$: $\Int A$ is open, $\closure{A}$ is closed, $\Int A$ is
	the largest open set contained in $A$, and $\closure A$ is the smallest closed
	set containing $A$.
\end{proposition}

\begin{proof}
	Let $a\in\Int A$, so $\ball{\varepsilon}{a}\subseteq A$ for some
	$\varepsilon>0$. By \autoref{prop:ball-open}, every $x\in\ball{\varepsilon}{a}$
	has a ball $\ball{\delta}{x}\subseteq\ball{\varepsilon}{a}\subseteq A$, so
	$x\in\Int A$; hence $\ball{\varepsilon}{a}\subseteq\Int A$ and $\Int A$ is open.
	Applying this to $\comp A$ and using \autoref{prop:interior-closure-duality}
	shows $\closure A$ is closed.

	If $G\subseteq A$ is open then every $a\in G$ has a ball inside $G\subseteq A$,
	so $G\subseteq\Int A$; this is maximality. The statement for $\closure A$
	follows by complementation.
\end{proof}

\begin{theorem}[The three axioms of a topology]\label{thm:topology-axioms}
	Let $\mathcal{O}$ denote the collection of open subsets of a metric space
	$(X,d)$. Then
	\begin{enumerate}[(i)]
		\item the union of any subfamily of $\mathcal{O}$ belongs to $\mathcal{O}$;
		\item the intersection of any \emph{finite} subfamily of $\mathcal{O}$ belongs
		      to $\mathcal{O}$;
		\item $\emptyset\in\mathcal{O}$ and $X\in\mathcal{O}$.
	\end{enumerate}
\end{theorem}

\begin{proof}
	(i) Let $(G_i)_{i\in I}\subseteq\mathcal O$ and $x\in G=\bigcup_iG_i$. Then
	$x\in G_{i_0}$ for some $i_0$, and openness of $G_{i_0}$ supplies
	$\varepsilon>0$ with $\ball{\varepsilon}{x}\subseteq G_{i_0}\subseteq G$.

	(ii) Let $G_1,\dots,G_n\in\mathcal O$ and $x\in G=\bigcap_{i=1}^nG_i$. For each
	$i$ pick $\varepsilon_i>0$ with $\ball{\varepsilon_i}{x}\subseteq G_i$ and set
	$\varepsilon=\min\{\varepsilon_1,\dots,\varepsilon_n\}$. Then $\varepsilon>0$
	--- this is where finiteness is used --- and
	$\ball{\varepsilon}{x}\subseteq\ball{\varepsilon_i}{x}\subseteq G_i$ for every
	$i$, so $\ball{\varepsilon}{x}\subseteq G$.

	(iii) For $X$: any $\varepsilon$ works since $\ball{\varepsilon}{x}\subseteq X$
	always. For $\emptyset$: the statement ``$x\in\emptyset$ implies
	$\ball{\varepsilon}{x}\subseteq\emptyset$'' is vacuously true, by the truth
	table of \autoref{sec:sets-logic}.
\end{proof}

\begin{eg}[Infinitely many open sets can intersect in a closed set]
	\label{eg:infinite-intersection}
	In $\R$, $\bigcap_{n\in\N}\left(-\tfrac1n,\tfrac1n\right)=\{0\}$, an
	intersection of open sets that is not open. The minimum in the proof of
	\autoref{thm:topology-axioms}(ii) becomes an infimum equal to $0$, and the
	argument collapses. Dually, $\bigcup_{n\in\N}\left[\tfrac1n,1\right]=(0,1]$ is a
	union of closed sets that is not closed.
\end{eg}

\begin{definition}[Topological space]\label{def:topological-space}
	Let $X$ be a set. A collection $\mathcal{O}\subseteq\pow{X}$ satisfying
	(i)--(iii) of \autoref{thm:topology-axioms} is a \dfn{topology} on $X$, the pair
	$(X,\mathcal{O})$ is a \dfn{topological space}, and the members of
	$\mathcal{O}$ are the open sets. A topology arising from a metric as above is
	\dfnas{metrisable topology}{metrisable}.
\end{definition}

\begin{remark}[Two abstractions of the same list of axioms]\label{rem:topology-sigma-algebra}
	Compare \autoref{def:topological-space} with the definition of a
	$\sigma$-algebra: a collection $\mathcal{F}\subseteq\pow{\Omega}$ containing
	$\emptyset$ and $\Omega$, closed under complementation, and closed under
	\emph{countable} unions. The two are different closure requirements on a
	collection of sets, and the difference is exactly where the two theories
	diverge: topology tolerates arbitrary unions but only finite intersections;
	measure theory tolerates countable unions and countable intersections but
	requires complements. On $\R$ the Borel $\sigma$-algebra is generated by the
	open sets, which is how the two are linked
	\autocite[slide 37]{fukuda2026analysis}. The parallel between the two
	definitions of ``measurable/continuous via preimages'' drawn there ---
	$f$ is continuous if $f^{-1}(B)$ is open whenever $B$ is open; $X$ is a random
	variable if $X^{-1}(B)$ is measurable whenever $B$ is measurable --- is made
	precise in \autoref{thm:continuity-preimage} and
	\autoref{def:random-variable}.
\end{remark}

\section{Sequences and Limits}\label{sec:top-sequences}

\begin{definition}[Sequence, subsequence, limit]\label{def:sequence}
	A \dfn{sequence} in $X$ is a map $\N\to X$, written $(x_n)_{n\in\N}$. If
	$\tau\colon\N\to\N$ is strictly increasing, then $(x_{\tau(n)})_{n\in\N}$ is a
	\dfn{subsequence}, usually written $(x_{n_k})_{k\in\N}$.

	The sequence \dfnas{convergent sequence}{converges} to $x^\ast\in X$, written
	$x_n\to x^\ast$ or $x^\ast=\lim_{n\to\infty}x_n$, if for every $\varepsilon>0$
	there is $n_0\in\N$ such that $n\geq n_0$ implies $d(x^\ast,x_n)<\varepsilon$.
\end{definition}

\begin{proposition}[Basic facts]\label{prop:sequence-basics}
	Let $(x_n)$ be a sequence in a metric space $X$.
	\begin{enumerate}[(i)]
		\item A limit, if it exists, is unique.
		\item If $x_n\to x^\ast$ then every subsequence converges to $x^\ast$.
		\item A convergent sequence is bounded.
	\end{enumerate}
\end{proposition}

\begin{proof}
	(i) If $x_n\to x$ and $x_n\to y$, then for every $\varepsilon>0$ we may choose
	$n$ with $d(x,x_n)<\varepsilon/2$ and $d(y,x_n)<\varepsilon/2$, so
	$d(x,y)\leq d(x,x_n)+d(x_n,y)<\varepsilon$. As $\varepsilon>0$ was arbitrary,
	$d(x,y)=0$, hence $x=y$.

	(ii) Since $\tau$ is strictly increasing, $\tau(k)\geq k$ by induction. Given
	$\varepsilon>0$ take $n_0$ from the definition; for $k\geq n_0$ we have
	$\tau(k)\geq k\geq n_0$, so $d(x^\ast,x_{\tau(k)})<\varepsilon$.

	(iii) Take $\varepsilon=1$ and the corresponding $n_0$. Then
	$d(x^\ast,x_n)<1$ for $n\geq n_0$, and setting
	$M\eqdef\max\{1,d(x^\ast,x_1),\dots,d(x^\ast,x_{n_0-1})\}$ gives
	$(x_n)\subseteq\ball{M+1}{x^\ast}$; the maximum is over a finite set and hence
	finite.
\end{proof}

Parts (ii) and (iii) are Exercises~7.30 and~7.31 of
\autocite[slide 43]{fukuda2026analysis}.

\begin{theorem}[Sequential characterisation of closedness]\label{thm:closed-sequential}
	A set $A\subseteq X$ is closed if and only if the limit of every convergent
	sequence of elements of $A$ lies in $A$.
\end{theorem}

\begin{proof}
	($\Rightarrow$) Let $A$ be closed, $(x_n)\subseteq A$ and $x_n\to x^\ast$. For
	every $\varepsilon>0$ there is $n$ with $x_n\in\ball{\varepsilon}{x^\ast}$, and
	$x_n\in A$, so $A\cap\ball{\varepsilon}{x^\ast}\neq\emptyset$. Hence
	$x^\ast\in\closure A=A$.

	($\Leftarrow$) Suppose the stated condition holds and let
	$x^\ast\in\closure A$. For each $n\in\N$ the set
	$A\cap\ball{1/n}{x^\ast}$ is non-empty; pick $x_n$ in it. Then
	$d(x^\ast,x_n)<1/n\to0$, so $x_n\to x^\ast$ with $(x_n)\subseteq A$, and the
	hypothesis gives $x^\ast\in A$. Hence $\closure A\subseteq A$, and the reverse
	inclusion always holds.
\end{proof}

This is Proposition~7.11 of \textcite{fukuda2026notes}
\autocite[slide 43]{fukuda2026analysis}. It converts a statement about balls into
a statement about sequences, and it is the form in which closedness is checked in
practice.

\subsection{Monotone subsequences and Bolzano--Weierstrass}

\begin{proposition}[Every real sequence has a monotone subsequence]\label{prop:monotone-subseq}
	Every sequence $(x_n)_{n\in\N}$ in $\R$ has a subsequence that is increasing or
	a subsequence that is decreasing (possibly both).
\end{proposition}

\begin{proof}
	Call $m\in\N$ a \emph{peak} if $x_m>x_n$ for all $n>m$, and let
	$S\eqdef\{m\in\N : m\text{ is a peak}\}$.

	If $S$ is infinite, enumerate it increasingly as $n_1<n_2<\cdots$. For each $k$,
	since $n_k$ is a peak and $n_{k+1}>n_k$, we get $x_{n_k}>x_{n_{k+1}}$; hence
	$(x_{n_k})$ is strictly decreasing.

	If $S$ is finite, let $N$ exceed every element of $S$ and set $n_1\eqdef N$.
	Since $n_1$ is not a peak, there is $n_2>n_1$ with $x_{n_2}\geq x_{n_1}$; since
	$n_2>N$ is not a peak either, there is $n_3>n_2$ with $x_{n_3}\geq x_{n_2}$, and
	so on. This produces an increasing subsequence.
\end{proof}

This is Proposition~7.12 of \textcite{fukuda2026notes}; the argument is the one
sketched on \autocite[slides 45--46]{fukuda2026analysis}, made explicit.

\begin{proposition}[Monotone convergence]\label{prop:monotone-convergence}
	An increasing sequence in $\R$ that is bounded above converges to
	$\sup_n x_n$; a decreasing sequence bounded below converges to $\inf_n x_n$.
\end{proposition}

\begin{proof}
	Let $(x_n)$ be increasing and bounded above and put $s\eqdef\sup_nx_n$, which
	exists by \autoref{ass:completeness-axiom}. Given $\varepsilon>0$,
	\autoref{prop:sup-char}(ii) supplies $n_0$ with $x_{n_0}>s-\varepsilon$. For
	$n\geq n_0$, monotonicity and the upper-bound property give
	$s-\varepsilon<x_{n_0}\leq x_n\leq s$, so $\abs{s-x_n}<\varepsilon$. The
	decreasing case follows by applying this to $(-x_n)$.
\end{proof}

\begin{theorem}[Bolzano--Weierstrass]\label{thm:bolzano-weierstrass}
	Every bounded sequence in $\R$ has a convergent subsequence. More generally,
	every bounded sequence in $\R^n$ has a convergent subsequence.
\end{theorem}

\begin{proof}
	In $\R$: by \autoref{prop:monotone-subseq} extract a monotone subsequence, which
	inherits boundedness, and apply \autoref{prop:monotone-convergence}.

	In $\R^n$: let $(x_m)_{m\in\N}$ be bounded, with coordinates
	$x_m=(x^1_m,\dots,x^n_m)$. Each coordinate sequence is bounded in $\R$, since
	$\abs{x^i_m}\leq\norm{x_m}$. Apply the one-dimensional case to
	$(x^1_m)_m$ to get a subsequence along which the first coordinate converges;
	apply it again, along that subsequence, to the second coordinate; and so on. As
	$n$ is finite, after $n$ extractions we obtain a single subsequence along which
	every coordinate converges. By \autoref{prop:coordinatewise} the subsequence
	converges in $\R^n$.
\end{proof}

\begin{proposition}[Coordinatewise convergence]\label{prop:coordinatewise}
	A sequence $(x_m)$ in $\R^n$ converges if and only if each coordinate sequence
	$(x^i_m)_m$ converges in $\R$, and then
	$\lim_mx_m=(\lim_mx^1_m,\dots,\lim_mx^n_m)$.
\end{proposition}

\begin{proof}
	If $x_m\to x^\ast$, then $\abs{x^i_m-x^{i\ast}}\leq\norm{x_m-x^\ast}\to0$ for
	each $i$. Conversely, if $x^i_m\to x^{i\ast}$ for each $i$, then given
	$\varepsilon>0$ choose $n_i$ with $\abs{x^i_m-x^{i\ast}}<\varepsilon/\sqrt n$
	for $m\geq n_i$; for $m\geq\max_in_i$,
	$\norm{x_m-x^\ast}=\bigl(\sum_i\abs{x^i_m-x^{i\ast}}^2\bigr)^{1/2}<\varepsilon$.
\end{proof}

This is Proposition~7.14 of \textcite{fukuda2026notes}
\autocite[slide 48]{fukuda2026analysis}. The finiteness of $n$ is essential: the
argument in \autoref{thm:bolzano-weierstrass} performs $n$ successive
extractions, and in an infinite-dimensional space it cannot be completed. Indeed
Bolzano--Weierstrass \emph{fails} there, as \autoref{eg:no-bw-infinite} shows.

\begin{eg}[Bolzano--Weierstrass fails in infinite dimensions]\label{eg:no-bw-infinite}
	In the space of square-summable sequences with
	$\norm{x}^2=\sum_ix_i^2$, consider the orthonormal vectors
	$e^1,e^2,\dots$, where $e^k$ has a $1$ in position $k$ and zeros elsewhere. Then
	$\norm{e^k}=1$ for all $k$, so the sequence is bounded, while
	$\norm{e^j-e^k}=\sqrt2$ for $j\neq k$, so no subsequence is Cauchy and hence
	none converges. Bounded and closed therefore does not imply compact outside
	finite dimensions, which is why \autoref{thm:heine-borel} is stated for $\R^n$
	and why the fixed-point arguments of \autoref{ch:dynamic} use completeness and
	contraction rather than compactness.
\end{eg}

\section{Compactness}\label{sec:top-compact}

\begin{definition}[Compactness]\label{def:compact}
	Let $(X,d)$ be a metric space and $K\subseteq X$.
	\begin{enumerate}[(i)]
		\item An \dfn{open cover} of $K$ is a family $(G_i)_{i\in I}$ of open sets with
		      $K\subseteq\bigcup_{i\in I}G_i$. The set $K$ is \dfn{compact} if every
		      open cover of $K$ admits a finite subcover, that is, a finite
		      $J\subseteq I$ with $K\subseteq\bigcup_{i\in J}G_i$.
		\item $K$ is \dfnas{sequentially compact set}{sequentially compact} if every
		      sequence in $K$ has a subsequence converging to a point of $K$.
	\end{enumerate}
\end{definition}

\begin{theorem}[Compactness and sequential compactness agree in metric spaces]
	\label{thm:compact-seq-compact}
	In a metric space, a set is compact if and only if it is sequentially compact.
\end{theorem}

The proof is standard and somewhat long; see
\textcite[Thm.~3.28 and its converse]{rudin1976principles} or
\textcite[Sec.~8.1]{fukuda2026notes}. In general \emph{topological} spaces the
two notions differ, which is one reason for staying within metric spaces
throughout this volume.

\begin{proposition}[Compact implies closed and bounded]\label{prop:compact-closed-bounded}
	In any metric space, a compact set is closed and bounded.
\end{proposition}

\begin{proof}
	Boundedness: fix $x_0\in X$; the balls $\ball{n}{x_0}$, $n\in\N$, cover $K$, and
	a finite subcover is contained in the largest of them.

	Closedness: let $x^\ast\in\closure K$ and take a sequence
	$(x_n)\subseteq K$ with $x_n\to x^\ast$, as in the proof of
	\autoref{thm:closed-sequential}. By sequential compactness
	(\autoref{thm:compact-seq-compact}) some subsequence converges to a point of
	$K$; by \autoref{prop:sequence-basics}(ii) that limit is $x^\ast$. Hence
	$x^\ast\in K$ and $K$ is closed by \autoref{thm:closed-sequential}.
\end{proof}

\begin{theorem}[Heine--Borel]\label{thm:heine-borel}
	A subset of $\R^n$ is compact if and only if it is closed and bounded.
\end{theorem}

\begin{proof}
	Necessity is \autoref{prop:compact-closed-bounded}. For sufficiency, let
	$K\subseteq\R^n$ be closed and bounded and let $(x_m)\subseteq K$. Boundedness
	and \autoref{thm:bolzano-weierstrass} give a convergent subsequence
	$x_{m_k}\to x^\ast$; closedness and \autoref{thm:closed-sequential} give
	$x^\ast\in K$. So $K$ is sequentially compact, hence compact by
	\autoref{thm:compact-seq-compact}.
\end{proof}

\begin{remark}[The equivalence is special to $\R^n$]\label{rem:heine-borel-fails}
	\autoref{eg:no-bw-infinite} exhibits a closed bounded set --- the closed unit
	ball of an infinite-dimensional space --- that is not compact. In fact the
	closed unit ball of a normed space is compact if and only if the space is
	finite-dimensional (Riesz's lemma; see
	\textcite[Ch.~5]{aliprantis2006infinite}). Whenever a proof in economics invokes
	``closed and bounded, hence compact'', the underlying space must be checked to
	be finite-dimensional.
\end{remark}

\begin{proposition}[Stability properties]\label{prop:compact-stability}
	\begin{enumerate}[(i)]
		\item A closed subset of a compact set is compact.
		\item A finite union of compact sets is compact; an arbitrary intersection of
		      compact sets is compact.
		\item A finite product of compact sets is compact in the product metric.
	\end{enumerate}
\end{proposition}

\begin{proof}
	(i) Let $F\subseteq K$ be closed with $K$ compact, and let $(x_n)\subseteq F$.
	Compactness of $K$ gives a subsequence converging to some $x^\ast\in K$;
	closedness of $F$ and \autoref{thm:closed-sequential} give $x^\ast\in F$.

	(ii) For a finite union, a sequence in $\bigcup_{j=1}^mK_j$ has infinitely many
	terms in some $K_{j_0}$, and that subsequence has a further subsequence
	converging in $K_{j_0}$. An intersection of compact sets is a closed subset
	(by \autoref{prop:compact-closed-bounded} and
	\autoref{prop:interior-closure-duality}) of any one of them, so (i) applies.

	(iii) Extract successively in each coordinate, as in the proof of
	\autoref{thm:bolzano-weierstrass}; finiteness of the number of factors is what
	makes the extraction terminate.
\end{proof}

\begin{eg}[The budget set is compact]\label{eg:budget-compact}
	For $p\in\R^n_{++}$ and $m>0$,
	$\budget(p,m)=\{x\in\R^n_+:\ip{p}{x}\leq m\}$ is compact. It is bounded, as
	shown after \autoref{prop:bounded-equiv}. It is closed: if
	$(x_k)\subseteq\budget(p,m)$ and $x_k\to x^\ast$, then
	$x^\ast\in\R^n_+$ because each coordinate is a limit of non-negative numbers,
	and $\ip{p}{x^\ast}=\lim_k\ip{p}{x_k}\leq m$ by continuity of the inner product
	(\autoref{prop:linear-continuous}). \autoref{thm:heine-borel} applies.

	Every hypothesis is doing work. If $p_i=0$ for some $i$ the set is unbounded. If
	the constraint were $\ip{p}{x}<m$ the set would not be closed and the supremum
	of a strictly increasing utility would not be attained. If consumption were not
	restricted to $\R^n_+$, the set would again be unbounded whenever some $p_i<0$.
	The existence half of \autoref{eg:consumer-setup} rests on this
	proposition, via the Weierstrass theorem of \autoref{ch:continuity}.
\end{eg}

\section{Connectedness}\label{sec:top-connected}

\begin{definition}[Connected set]\label{def:connected}
	A metric space $X$ is \dfn{connected} if the only subsets of $X$ that are both
	open and closed are $\emptyset$ and $X$; equivalently, $X$ cannot be written as
	the union of two disjoint non-empty open sets. A subset $A\subseteq X$ is
	connected if it is connected as a metric space in the induced metric.
\end{definition}

\begin{proposition}\label{prop:connected-interval}
	A subset of $\R$ is connected if and only if it is an interval, that is, if
	$x,y\in A$ and $x<z<y$ imply $z\in A$. Every convex subset of $\R^n$ is
	connected.
\end{proposition}

\begin{proof}
	\emph{Connected $\Rightarrow$ interval.} If $A$ is not an interval, take
	$x<z<y$ with $x,y\in A$ and $z\notin A$; then $A\cap(-\infty,z)$ and
	$A\cap(z,\infty)$ are disjoint, non-empty, relatively open, and cover $A$, so
	$A$ is not connected.

	\emph{Interval $\Rightarrow$ connected.} This direction is proved from the
	supremum property of $\R$ (\autoref{ass:completeness-axiom}), not from the
	connectedness of $[0,1]$, which is the same statement. Let $A$ be an interval
	and suppose $A=U\cup W$ with $U,W$ relatively open, disjoint and non-empty. Pick
	$x\in U$ and $y\in W$; without loss of generality $x<y$, and $[x,y]\subseteq A$
	because $A$ is an interval. Put
	\[
		s\eqdef\sup\bigl\{t\in[x,y] : [x,t]\subseteq U\bigr\},
	\]
	the set being non-empty, since $x$ belongs to it, and bounded above by $y$; so
	$s$ is a real number lying in $[x,y]\subseteq A$. Since $A=U\cup W$, either
	$s\in W$ or $s\in U$.

	Suppose $s\in W$. Relative openness gives $\varepsilon>0$ with
	$A\cap(s-\varepsilon,s+\varepsilon)\subseteq W$, and the definition of the
	supremum gives $t\in[x,y]$ with $s-\varepsilon<t\leq s$ and $[x,t]\subseteq U$.
	Then $t\in A\cap(s-\varepsilon,s+\varepsilon)\subseteq W$ and $t\in U$, so
	$t\in U\cap W=\emptyset$.

	Suppose instead $s\in U$, and take $\varepsilon>0$ with
	$A\cap(s-\varepsilon,s+\varepsilon)\subseteq U$ and $t$ as above. Put
	$s'\eqdef\min\{s+\varepsilon/2,\ y\}$. Every point of $[t,s']$ lies in
	$[x,y]\subseteq A$ and in $(s-\varepsilon,s+\varepsilon)$, hence in $U$;
	together with $[x,t]\subseteq U$ this gives $[x,s']\subseteq U$ and therefore
	$s'\leq s$. If $s<y$ then $s'>s$, a contradiction; so $s=y$ and $y\in U$,
	contradicting $y\in W$.

	Both cases are impossible, so no such decomposition exists and $A$ is
	connected.

	\emph{Convex sets in $\R^n$.} Let $C\subseteq\R^n$ be convex and non-empty and
	suppose $C=U\cup W$ with $U,W$ relatively open, disjoint and non-empty. Pick
	$x\in U$, $y\in W$ and consider $\gamma(t)\eqdef(1-t)x+ty$, which is continuous
	and maps $[0,1]$ into $C$. Then $\gamma^{-1}(U)$ and $\gamma^{-1}(W)$ are
	relatively open in $[0,1]$ by \autoref{thm:continuity-preimage}, disjoint,
	non-empty and cover $[0,1]$, contradicting connectedness of the interval
	$[0,1]$ just established.
\end{proof}

Connectedness is what makes the intermediate value theorem true, and hence what
underlies both the existence of an interior solution between two boundary
regimes and Debreu's continuous utility representation theorem alluded to in
\autoref{rem:lexicographic}.

\section{Chapter Summary}\label{sec:top-summary}

Open sets are those equal to their interior, closed sets those equal to their
closure, and the two notions are exchanged by complementation. The open sets are
closed under arbitrary unions and finite intersections and contain
$\emptyset$ and $X$; abstracting these three properties gives a topology, and the
countable-union-plus-complement variant gives a $\sigma$-algebra. A set is closed
if and only if it contains the limits of its convergent sequences. Every bounded
real sequence has a monotone, hence convergent, subsequence
(\autoref{thm:bolzano-weierstrass}); by coordinatewise extraction the same holds
in $\R^n$, and the argument requires $n$ to be finite. Compactness is equivalent
to sequential compactness in metric spaces, implies closed and bounded always,
and is equivalent to closed and bounded exactly in $\R^n$
(\autoref{thm:heine-borel}). The budget set with strictly positive prices is
compact, which is the hypothesis under which \autoref{ch:continuity} produces an
optimal consumption bundle.

\begin{exercise}\label{ex:top-1}
	Show that $\partial A=\closure{A}\cap\closure{\comp A}$, that
	$\partial A$ is always closed, and that $A$ is closed if and only if
	$\partial A\subseteq A$. Compute $\Int A$, $\closure A$ and $\partial A$ for
	$A=\Q\cap[0,1]$ in $\R$, and comment on which of the three is empty.
\end{exercise}

\begin{exercise}\label{ex:top-2}
	Prove that in the discrete metric of \autoref{eg:discrete-metric} every subset
	is both open and closed, that a sequence converges if and only if it is
	eventually constant, and that a set is compact if and only if it is finite.
	Which hypothesis of \autoref{thm:heine-borel} fails for an infinite bounded set
	in this metric?
\end{exercise}

\begin{exercise}\label{ex:top-3}
	Let $K\subseteq\R^n$ be compact and non-empty. Show that $\sup K$ and $\inf K$
	are attained coordinatewise, that is, for each $i$ there exist
	$x,y\in K$ with $x_i=\max\{z_i : z\in K\}$ and $y_i=\min\{z_i: z\in K\}$. Give an
	example of a closed but unbounded $K\subseteq\R$ for which the conclusion fails.
\end{exercise}
